% =====================================================================================================
\section{Introduction}\label{sec:intro}

Bayesian variable selection in the linear model with many more predictors than observations has a large theory, but
that theory is mostly about an exact posterior under a prior the analyst tunes. Users run software with fixed defaults,
and read what it reports: inclusion probabilities, a median probability model, an evidence estimate, a best model. This
paper proves consistency results for one method exactly as its software implements it.

\paragraph{The method.} SMC-BVS puts a fixed prior on models and coefficients, a beta-binomial prior on the model size
\citep{ScottBerger2010} and a $g$-prior with a unit ridge, so that the posterior over models is known in closed form up to
its normalising constant, the evidence. Nothing is tuned by the user. The method partitions the predictors into small
blocks of correlated predictors, scores every subset of every block exactly, and combines these block tables into a
\emph{start law}, an approximation of the posterior that can be computed and sampled exactly. A sequential Monte Carlo
sampler, run as four independent \emph{islands}, then moves from the start law to the posterior along a tempered path
that an internal calibration fixes in advance. Two deterministic local searches, one of them from the mode of the start
law, finish the run. The run reports its best model, pooled inclusion probabilities, their median probability model and
an evidence estimate (Section~\ref{sec:method}).

\paragraph{What we prove.} Under conditions stated in Section~\ref{sec:results}:
\begin{enumerate}[leftmargin=1.6em,itemsep=1pt]
  \item the posterior puts all but a fraction of order $\rn$ of its mass on the true model, and its highest posterior and
    median probability models are the true model with probability tending to one, on every design of our simulation
    studies, strong common factors included (Result~1);
  \item the posterior recovers the true conditional law of the response at radius of order $\rn$ in Hellinger distance,
    with no lower bound on the coefficients (Result~2);
  \item with probability tending to one, every run that produces output reports exactly the true model, for every seed
    and every outcome of its calibration, when the predictors that can imitate the signal are few and local; pervasive
    common factors with $p\gg\sqrt n$ are excluded (Result~3);
  \item the evidence estimate is unbiased on every design, and the reported inclusion probabilities, median model and
    evidence estimate are consistent on designs without leakage, on the event that the calibration succeeds (Result~4);
  \item no design allows a posterior rate faster than $n^{-1/2}$, and on two families of designs the rate $\rn$ of
    Result~1 is attained (Result~5).
\end{enumerate}

\paragraph{Earlier work.} The closest earlier analysis is that of \citet{SongLiang2015}, who prove selection and
distribution consistency of a posterior with $p\gg n$ and run a sampler on it. Their statements about what their
sampler reports are limits as the number of iterations grows, at fixed data and for every data set; ours hold at the
budget the software actually uses, jointly with $n$, under conditions on the design. Neither kind of statement implies
the other. Our results concern a different posterior (the prior as the software sets it), and Section~\ref{sec:SL-concl}
compares the conclusions theorem by theorem. [A paragraph on the wider literature on selection consistency, on Markov
chain Monte Carlo for variable selection and on sequential Monte Carlo over binary spaces awaits the owner's reference
list.]

\paragraph{Are the assumptions reasonable?} Our theorems use the same kinds of assumptions as \citet{SongLiang2015}; on
the designs we study, the assumptions of the posterior theorems hold, and where the two papers differ the direction
depends on the regime (Section~\ref{sec:assumptions}). The designs are those of our simulation studies, a study of
$2{,}400$ fits, a study of strong common factors, and a planned set of runs with $p$ up to $10^6$, which are reported
separately; this paper contains no simulation study.

\paragraph{Organisation.} Section~\ref{sec:setup} sets up the data and the notation, and Section~\ref{sec:method}
describes the method as coded. Section~\ref{sec:results} states the five results, Section~\ref{sec:assumptions} judges
their assumptions, and Section~\ref{sec:scope} states what is not claimed. The supplement (Appendix~\ref{sec:supp})
holds further results, and the appendices hold the notation, the full parameter table with its code locations, the
explicit constants and the status of each result.
